\chapter{Аудит космологической постоянной как якоря проводимости}
\label{ch:v10-cosmological-conductance-anchor-audit}

\section{Точная космологическая формула ветви}

Общий родитель предыдущего гейта выбрал
\begin{equation}
 \kappa=\Gamma_B=3\Delta\zeta\,\Omega,
 \qquad H_B=\Delta\zeta\,\Omega=\frac{\kappa}{3}.
 \label{eq:v10-lambda-anchor-rate-chain}
\end{equation}
Если космологическую постоянную растущей ветви определить стандартной
формулой кривизны
\begin{equation}
 \Lambda_{\rm growth}=3\left(\frac{H_B}{c}\right)^2,
 \label{eq:v10-lambda-anchor-definition}
\end{equation}
то немедленно получаем
\begin{equation}
 \Lambda_{\rm growth}=\frac{\kappa^2}{3c^2},
 \qquad
 c\sqrt{3\Lambda_{\rm growth}}=\kappa.
 \label{eq:v10-lambda-anchor-tautology}
\end{equation}
Совпадение размерности и коэффициента точное. Но второе равенство лишь
обращает первое: $\Lambda_{\rm growth}$ вычислена из той же проводимости,
которую предполагается восстановить.

Соответствующий радиус
\begin{equation}
 \ell_\Lambda:=\sqrt{\frac{3}{\Lambda_{\rm growth}}}
\end{equation}
удовлетворяет
\begin{equation}
 \kappa\ell_\Lambda=3c.
 \label{eq:v10-lambda-anchor-radius-product}
\end{equation}
Следовательно, без независимо выбранного радиуса фиксируется только
произведение проводимости и длины.

\section{Аудит семи космологических кандидатов}

Проверены семь вариантов: космологическая постоянная роста, собственная
кривизна ячейки, спектральная кривизна обрезания, кривизна индуцированного
сквозным током вакуума, наблюдаемая космологическая постоянная,
планковско-вакуумная кривизна и топологическая плотность кривизны.

Каждый кандидат оценён по шести условиям:
\begin{enumerate}
\item правильная размерность кривизны;
\item наличие внутри построенного носителя;
\item независимость от $\kappa$, $\Gamma_B$ и $\Omega$;
\item типизированное отображение в проводимость;
\item разрыв общей частотной орбиты;
\item происхождение из общего родителя без внешней цели.
\end{enumerate}

Матрица условий имеет размер $7\times6$ и ранг $5$. Полных проходов нет;
максимальный результат равен $4/6$. Наблюдаемая $\Lambda$ и планковская
кривизна действительно разорвали бы орбиту, но являются внешними входами.
Внутренние кандидаты либо зависят от уже выбранной скорости, либо содержат
собственный свободный радиус, обрезание или квант объёма.

\section{Масштабный ранг}

В логарифмических координатах
$(\kappa,\Gamma_B,\Omega,H_B,\Lambda)$ установленные отношения имеют матрицу
\begin{equation}
 C_\Lambda=
 \begin{pmatrix}
 1&-1&0&0&0\\
 0&1&-1&0&0\\
 0&0&-1&1&0\\
 0&0&0&-2&1
 \end{pmatrix}.
 \label{eq:v10-lambda-anchor-scale-map}
\end{equation}
Она имеет ранг $4$, ядро размерности $1$ и уничтожает вектор
\begin{equation}
 (1,1,1,1,2)^T.
 \label{eq:v10-lambda-anchor-scale-vector}
\end{equation}
То есть все скорости масштабируются как $c$, а $\Lambda$ --- как $c^2$.
Добавление независимого условия, фиксирующего $\Lambda$, повышает ранг до
$5$ и действительно выбирает абсолютную проводимость. Необходимая форма
якоря установлена, но его внутреннее происхождение не найдено.

\begin{theorem}[Условная космологическая калибровка проводимости]
Если положительная космологическая постоянная выведена независимо от часов,
рождения ячеек и хопфовского тока, то она фиксирует
$\kappa=c\sqrt{3\Lambda}$. Дополнительное абсолютное условие устраняет
последнюю масштабную моду.
\end{theorem}

\begin{nogocriterion}[Запрет обратного чтения собственной кривизны роста]
Величина $\Lambda_{\rm growth}=\kappa^2/(3c^2)$ не может служить
независимым объяснением $\kappa$: равенство
$c\sqrt{3\Lambda_{\rm growth}}=\kappa$ является тождественным обращением
определения. Геометрический образ становится физическим якорем только после
вывода $\Lambda$ из другого сектора общего родителя.
\end{nogocriterion}

\begin{protocol}[Статус космологического якоря]\begin{enumerate}
\item проверенных кандидатов: \textbf{$7$};
\item полных проходов: \textbf{$0/7$};
\item максимальный результат: \textbf{$4/6$};
\item ранг/ядро относительной карты: \textbf{$4/1$};
\item условный внешний якорь $\Lambda$: \textbf{$1/1$};
\item физическое происхождение $\Lambda$: \textbf{$0/1$};
\item абсолютная проводимость: \textbf{$0/1$};
\item следующий гейт: \textbf{родитель космологической постоянной из
сквозного потока}.
\end{enumerate}\end{protocol}

Машинный сертификат сохранён в
\path{s2t_v10_cell_birth_four_volume_cosmological_constant_conductance_anchor_candidate_audit_gate_results.json}.